% This file was assembled from main.tex, custom.sty, manuscript inputs, and bib/references.bib.
% Generated on 2026-06-14.

% The filecontents* block writes the embedded bibliography to references.bib.
% This lets the exported tex file compile without the original bib/references.bib file.
\begin{filecontents*}{references.bib}
@misc{ABGLSSTWM26,
  author       = {Noga Alon and Thomas F. Bloom and W. T. Gowers and Daniel Litt and Will Sawin and Arul Shankar and Jacob Tsimerman and Victor Wang and Melanie Matchett Wood},
  title        = {Remarks on the disproof of the unit distance conjecture},
  year         = {2026},
  eprint       = {2605.20695},
  archivePrefix = {arXiv}
}

@misc{Bloom26,
  author       = {T. F. Bloom},
  title        = {{Erd{\H{o}}s Problem {\#}1208}},
  year         = {2026},
  howpublished = {\url{https://www.erdosproblems.com/1208}},
  note         = {Accessed 2026-06-14}
}

@article{Erdos46,
  author  = {Paul Erd\H{o}s},
  title   = {On sets of distances of {$n$} points},
  journal = {American Mathematical Monthly},
  volume  = {53},
  number  = {5},
  pages   = {248--250},
  year    = {1946}
}

@article{ErGu70,
  author  = {Paul Erd\H{o}s and R. K. Guy},
  title   = {Distinct distances between lattice points},
  journal = {Elemente der Mathematik},
  volume  = {25},
  pages   = {121--123},
  year    = {1970}
}

@article{LeTh95,
  author  = {Hanno Lefmann and Torsten Thiele},
  title   = {Point sets with distinct distances},
  journal = {Combinatorica},
  volume  = {15},
  number  = {3},
  pages   = {379--408},
  year    = {1995},
  doi     = {10.1007/BF01299744}
}

@article{Ch13,
  author  = {Marcos Charalambides},
  title   = {A note on distinct distance subsets},
  journal = {Journal of Geometry},
  volume  = {104},
  number  = {3},
  pages   = {439--442},
  year    = {2013},
  doi     = {10.1007/s00022-013-0176-0},
  eprint  = {1211.1776},
  archivePrefix = {arXiv}
}

@article{CFGHUZ15,
  author  = {David Conlon and Jacob Fox and William Gasarch and David G. Harris and Douglas Ulrich and Samuel Zbarsky},
  title   = {Distinct volume subsets},
  journal = {SIAM Journal on Discrete Mathematics},
  volume  = {29},
  number  = {1},
  pages   = {472--480},
  year    = {2015}
}

@article{GuthKatz15,
  author  = {Larry Guth and Nets Hawk Katz},
  title   = {On the {E}rd{\H{o}}s distinct distances problem in the plane},
  journal = {Annals of Mathematics},
  volume  = {181},
  number  = {1},
  pages   = {155--190},
  year    = {2015}
}

@misc{CFRN26,
  author       = {Felix Christian Clemen and Jakob F\"uhrer and Oliver Roche-Newton},
  title        = {Geometric {S}idon Problems},
  year         = {2026},
  eprint       = {2606.05841},
  archivePrefix = {arXiv}
}

@misc{OpenAIUnitDistance26,
  author = {{OpenAI}},
  title  = {Planar Point Sets with Many Unit Distances},
  year   = {2026},
  note   = {Available from OpenAI}
}

@misc{PohoataShefferBlog26,
  author = {Cosmin Pohoata and Adam Sheffer},
  title  = {Sets with no isosceles triangles / repeated distances in Atlanta},
  year   = {2026},
  note   = {Blog post, posted April 27, 2026 and updated June 11, 2026}
}

@misc{CMPShY26,
  author = {Ernie Croot and Junzhe Mao and Cosmin Pohoata and Adam Sheffer and Chi Hoi Yip},
  title  = {A combinatorial large sieve for {S}idon sets, distances, and norm forms},
  year   = {2026},
  note   = {Preprint}
}

@misc{Pohoata26SplitPrimes,
  author       = {Cosmin Pohoata},
  title        = {Split primes and the {E}lekes--{R}\'onyai problem},
  year         = {2026},
  eprint       = {2606.13619},
  archivePrefix = {arXiv}
}

@article{HMR21,
  author  = {Farshid Hajir and Christian Maire and Ravi Ramakrishna},
  title   = {Cutting towers of number fields},
  journal = {Annales Mathématiques du Québec},
  volume  = {45},
  pages   = {321--345},
  year    = {2021}
}

@book{Neu99,
  author    = {J{\"u}rgen Neukirch},
  title     = {Algebraic Number Theory},
  series    = {Grundlehren der mathematischen Wissenschaften},
  volume    = {322},
  publisher = {Springer-Verlag},
  address   = {Berlin},
  year      = {1999}
}

@misc{Mil20,
  author = {J. S. Milne},
  title  = {Algebraic Number Theory},
  year   = {2020},
  note   = {Version 3.08, available at \url{https://www.jmilne.org/math/CourseNotes/ANT.pdf}}
}

@book{RZ10,
  author    = {Luis Ribes and Pavel Zalesskii},
  title     = {Profinite Groups},
  edition   = {2},
  series    = {Ergebnisse der Mathematik und ihrer Grenzgebiete. 3. Folge},
  volume    = {40},
  publisher = {Springer-Verlag},
  address   = {Berlin},
  year      = {2010}
}
\end{filecontents*}

\documentclass{amsart}


% ====================
% Package loading
% ====================

% Mathematics
\usepackage{amsmath, amssymb, amsthm}

% Hyperlinks
\usepackage[colorlinks=true,citecolor=blue]{hyperref}

% ====================
% Configuration
% ====================

% Display math settings
\allowdisplaybreaks

% Theorem environments
\theoremstyle{plain}
\newtheorem{theorem}{Theorem}[section]
\newtheorem{lemma}[theorem]{Lemma}
\newtheorem{proposition}[theorem]{Proposition}
\newtheorem{corollary}[theorem]{Corollary}
\newtheorem*{claim}{Claim}

\theoremstyle{definition}
\newtheorem{definition}[theorem]{Definition}
\newtheorem{example}[theorem]{Example}

\theoremstyle{remark}
\newtheorem{remark}[theorem]{Remark}
\newtheorem*{note}{Note}
% \usepackage{draftview} % Comment out this line for the final view.

\title{Planar Point Sets with Small Distinct-Distance Subsets}
\author{Sungchul Lee}
\date{\today}

\begin{document}

\begin{abstract}
Let \(F_2(n)\) denote the minimum, over all \(n\)-point sets in the plane, of the largest size of a subset whose pairwise distances are all distinct.
The classical square-grid construction gives \(F_2(n)\ll n^{1/2}/(\log n)^{1/4}\).
We prove that this can be improved by a fixed power: there is an absolute constant \(\varepsilon>0\) such that \(F_2(n)\ll n^{1/2-\varepsilon}\).
The proof combines two ingredients: the number-field towers of bounded root discriminant used in the OpenAI unit-distance construction, and an argument of Pohoata and Sheffer based on split primes.
\end{abstract}

\maketitle

% \tableofcontents

\section{Introduction}
\label{sec:introduction}

For \(n\ge1\), let \(F_2(n)\) be the minimum, over all \(n\)-point sets \(P\subset\mathbb R^2\), of the maximum size of a subset of \(P\) whose pairwise distances are all distinct.
We call such a subset \emph{distance-Sidon}.
Equivalently, every \(n\)-point set in the plane contains a distance-Sidon subset of size at least \(F_2(n)\).
The problem is to determine the order of magnitude of \(F_2(n)\) as \(n\to\infty\).
This is the planar case of Erd\H{o}s' distance-Sidon subset problem \cite{Bloom26}, rooted in the classical study of distances determined by finite point sets \cite{Erdos46}, \cite{ErGu70}; see also \cite{LeTh95} and \cite{CFGHUZ15}.

Until recently, the long-standing bounds were
\[
\frac{n^{1/3}}{(\log n)^{1/3}}
\ll F_2(n)
\ll
\frac{n^{1/2}}{(\log n)^{1/4}}.
\]
The upper bound comes from the square integer grid: a \(\sqrt n\times\sqrt n\) grid determines only \(O(n/\sqrt{\log n})\) squared distances, by the classical estimate for integers representable as a sum of two squares, and hence no distance-Sidon subset of the grid can have more than \(O(n^{1/2}/(\log n)^{1/4})\) points.
The lower bound was proved by Charalambides \cite{Ch13}, using the Guth--Katz distinct-distance machinery \cite{GuthKatz15}.
Recently Clemen, F\"uhrer and Roche-Newton \cite{CFRN26} removed the logarithmic loss on the lower-bound side, proving \(F_2(n)\gg n^{1/3}\).

Our main result is a polynomial improvement over the upper bound.

\begin{theorem}\label{thm:main}
There exists an absolute constant \(\varepsilon>0\) such that
\[
F_2(n)\ll n^{1/2-\varepsilon}.
\]
\end{theorem}

The proof combines two recent ideas.
The first is the number-field framework used in the OpenAI construction disproving Erd\H{o}s' unit-distance conjecture \cite{OpenAIUnitDistance26}, \cite{ABGLSSTWM26}.
It provides totally real fields \(L\) of growing degree and bounded root discriminant, together with rational primes \(q\equiv1\pmod4\) which split completely in \(L\).
We use these fields and split primes to construct point sets with small distance-Sidon subsets.

The second idea is the split-prime isotropic-line sieve of Pohoata and Sheffer for the integer grid \cite{PohoataShefferBlog26}.
In their argument one chooses primes \(p_i\equiv1\pmod4\), roots \(x_i^2\equiv -1\pmod {p_i}\), and for each sign vector \(\sigma\in\{\pm1\}^t\) considers the congruence system
\[
u\equiv \sigma_i x_i v\pmod {p_i}\qquad(1\le i\le t).
\]
These are isotropic directions modulo \(d=\prod_i p_i\), and a double count over the \(2^t\) sign vectors bounds distance-Sidon subsets of \([N]^2\) by
\[
N\exp\left(-c\frac{\log N}{\log\log N}\right)
\]
for an absolute constant \(c>0\).
See also the combinatorial large-sieve framework announced in \cite{CMPShY26} and the related split-prime amplification in \cite{Pohoata26SplitPrimes}.

Our argument lifts this sieve from \(\mathbb Z^2\) to \(\mathcal O_L^2\).
If \(q_i\) splits completely in \(L\), then it gives \([L:\mathbb Q]\) prime ideals, and therefore the \(2^t\) local sign choices in the grid become \(2^{t[L:\mathbb Q]}\) choices.
This tower amplification is the main point: a fixed set of split primes yields a saving exponential in the field degree, and that exponential-in-degree saving becomes a fixed power saving in the number of planar points.

\subsection*{Organization}
Section~\ref{sec:preliminaries} records the required facts about the number-field tower.
Section~\ref{sec:point-set-construction} constructs the planar point sets from algebraic boxes.
Section~\ref{sec:sieve-estimate} proves the tower-amplified split-prime sieve.
Section~\ref{sec:proof-main-theorem} chooses the parameters and proves Theorem~\ref{thm:main}.

\subsection*{Notation}
For a finite set \(S\), we write \(\#S\) for its cardinality.

\section{Preliminaries}
\label{sec:preliminaries}

\subsection{Basic algebraic number theory}
\label{subsec:basic-algebraic-number-theory}

We use standard terminology from algebraic number theory; see, for example, \cite[Chapter~I]{Neu99} or \cite{Mil20}.

A \emph{number field} \(F\) is a finite extension of \(\mathbb Q\).
Put \(d=[F:\mathbb Q]\), and let \(\sigma_1,\ldots,\sigma_d:F\hookrightarrow\mathbb C\) be the embeddings of \(F\) into \(\mathbb C\).
For \(\alpha\in F\), the \emph{trace} of \(\alpha\) is
\[
\operatorname{Tr}_{F/\mathbb Q}(\alpha)
=
\sum_{k=1}^d\sigma_k(\alpha).
\]
Let \(\mathcal O_F\) be the ring of integers of \(F\), i.e.,
\[
\mathcal O_F
=
\{x\in F:x\text{ is a zero of some monic polynomial over }\mathbb Z\}.
\]
Then \(\mathcal O_F\) has a \(\mathbb Z\)-basis \(\omega_1,\ldots,\omega_d\).
The \emph{discriminant} of \(F\) is
\[
\Delta_F
=
\det\bigl(\operatorname{Tr}_{F/\mathbb Q}(\omega_i\omega_j)\bigr)_{1\le i,j\le d}.
\]
This integer is independent of the choice of integral basis.
The \emph{root discriminant} of \(F\) is
\[
|\Delta_F|^{1/[F:\mathbb Q]}.
\]
A number field \(F\) is said to be \emph{totally real} if every embedding \(F\hookrightarrow\mathbb C\) has image in \(\mathbb R\).

A \emph{rational prime} is an ordinary prime number in \(\mathbb Z\).
Let \(F/\mathbb Q\) be a number field of degree \(d\), and let \(q\) be a rational prime.
The ring \(\mathcal O_F\) is a Dedekind domain, so every nonzero ideal factors uniquely as a product of prime ideals, and every nonzero prime ideal is maximal; see, for example, \cite[Theorem~3.29, Theorem~3.7, and Definition~3.3]{Mil20}.
We say that \(q\) \emph{splits completely} in \(F\) if
\[
q\mathcal O_F=\mathfrak q_1\cdots\mathfrak q_d
\]
as a product of \(d\) distinct prime ideals of \(\mathcal O_F\).
In this case \(\mathcal O_F/\mathfrak q_k\cong\mathbb F_q\) for \(k=1,\ldots,d\).\footnote{Indeed, if \(q\mathcal O_F=\prod_k\mathfrak q_k^{e_k}\) and \(\#(\mathcal O_F/\mathfrak q_k)=q^{f_k}\), then \(\#(\mathcal O_F/q\mathcal O_F)=q^d\) and the Chinese remainder theorem gives \(d=\sum_k e_kf_k\). In the completely split case there are \(d\) distinct prime ideals and each \(e_k=1\), so every \(f_k=1\).}
For an infinite algebraic extension \(M/\mathbb Q\), the prime \(q\) is said to split completely in \(M\) if it splits completely in every finite subextension of \(M\).

\subsection{Pro-\texorpdfstring{\(p\)}{p} extensions}
\label{subsec:pro-p-extensions}

We use standard terminology for pro-\(p\) groups; see, for example, \cite[Chapters~1, 2, and~4]{RZ10}.

An inverse system consists of groups \(G_i\) and homomorphisms \(\phi_{ij}:G_j\to G_i\) for \(i\le j\), satisfying \(\phi_{ii}=\mathrm{id}\) and \(\phi_{ik}=\phi_{ij}\circ\phi_{jk}\) for \(i\le j\le k\).
Its inverse limit is the group of tuples \((g_i)_i\in\prod_iG_i\) such that \(\phi_{ij}(g_j)=g_i\) for all \(i\le j\).
A \emph{pro-\(p\) group} is an inverse limit of finite \(p\)-groups, where each finite \(p\)-group is given the discrete topology and the inverse limit is given the subspace topology inherited from \(\prod_iG_i\).
A \emph{Galois pro-\(p\) extension} \(M/K\) is a Galois extension whose Galois group is a pro-\(p\) group.

\subsection{A tower of totally real fields}
\label{subsec:number-field-tower}

\cite[Proposition~2.3]{ABGLSSTWM26} states the following result.

\begin{proposition}\label{prop:abglsstwm-tower}
There exists an infinite tower of totally real fields over \(\mathbb Q\) with bounded root discriminant and with infinitely many completely split primes \(q\equiv1\pmod4\).
\end{proposition}

Although the following proposition is not stated separately in \cite{ABGLSSTWM26}, it follows because the proof of \cite[Proposition~2.3]{ABGLSSTWM26} obtains the tower from an infinite Galois pro-\(2\) extension of \(\mathbb Q\), constructed using \cite[Theorem~2.8 of the arXiv version, or Theorem~4 of the published version]{HMR21}.

\begin{proposition}\label{prop:tower}
There exists an infinite tower \(\mathbb Q=F_0\subset F_1\subset F_2\subset\cdots\) of totally real fields over \(\mathbb Q\) such that the following hold.
\begin{enumerate}
\item The tower has bounded root discriminant.
\item Infinitely many rational primes \(q\equiv1\pmod4\) split completely in every \(F_j\).
\item \([F_j:\mathbb Q]=2^j\) for every \(j\ge0\).
\end{enumerate}
\end{proposition}

\begin{proof}
The proof of Proposition~\ref{prop:abglsstwm-tower} given in \cite{ABGLSSTWM26} supplies an infinite Galois pro-\(2\) extension \(M/\mathbb Q\) whose intermediate fields \(E\) with \(\mathbb Q\subset E\subset M\) and \([E:\mathbb Q]<\infty\) form a tower of totally real fields with bounded root discriminant and with infinitely many completely split primes \(q\equiv1\pmod4\).

It remains only to choose finite layers with degree exactly \(2^j\).
Let \(G=\operatorname{Gal}(M/\mathbb Q)\).
Set \(G_0=G\), and suppose \(G_j\) has been chosen.
Since \(G_j\) is an infinite pro-\(2\) group, it has an open subgroup \(G_{j+1}\subset G_j\) of index \(2\): take a nontrivial finite \(2\)-group quotient of \(G_j\), then take the inverse image of an index-\(2\) subgroup.
Put \(F_j=M^{G_j}\).\footnote{Here \(M^{G_j}=\{x\in M:\sigma(x)=x\text{ for every }\sigma\in G_j\}\).}
Then
\[
[F_{j+1}:F_j]=[G_j:G_{j+1}]=2,
\]
so \([F_j:\mathbb Q]=2^j\).
\end{proof}

\section{The point set construction}
\label{sec:point-set-construction}

We fix the tower \(\mathbb Q=F_0\subset F_1\subset F_2\subset\cdots\) from Proposition~\ref{prop:tower}.
Choose \(D\ge1\) such that every field in this tower has root discriminant at most \(D\).

\subsection{A translated algebraic box}
\label{subsec:translated-box}

Let \(j\ge0\).
Write the real embeddings of \(F_j\) as
\begin{equation}\label{eq:real-embeddings}
\sigma_{j,1},\ldots,\sigma_{j,2^j}:F_j\hookrightarrow\mathbb R
\end{equation}
and let
\begin{equation}\label{eq:sigma-map}
\Sigma_j:F_j\to\mathbb R^{2^j},
\qquad
\Sigma_j(\alpha)=(\sigma_{j,1}(\alpha),\ldots,\sigma_{j,2^j}(\alpha)).
\end{equation}
Then \(\Sigma_j(\mathcal O_{F_j})\) is a lattice of covolume\footnote{Covolume is the volume of a fundamental domain. Since \(F_j\) is totally real, if \(\omega_1,\ldots,\omega_{2^j}\) is a \(\mathbb Z\)-basis of \(\mathcal O_{F_j}\) and \(A=(\sigma_{j,k}(\omega_\ell))_{1\le k,\ell\le 2^j}\), then \(A\) is the real coordinate matrix of the lattice basis \(\Sigma_j(\omega_1),\ldots,\Sigma_j(\omega_{2^j})\), so \(\operatorname{covol}(\Sigma_j(\mathcal O_{F_j}))=|\det A|\). Also \(\bigl(\operatorname{Tr}_{F_j/\mathbb Q}(\omega_k\omega_\ell)\bigr)_{1\le k,\ell\le 2^j}=A^TA\). Hence \(|\Delta_{F_j}|=|\det(A^TA)|=|\det A|^2\).}
\begin{equation}\label{eq:covolume-bound}
\operatorname{covol}(\Sigma_j(\mathcal O_{F_j}))
=|\Delta_{F_j}|^{1/2}
\le D^{2^j/2}.
\end{equation}

Fix \(R\ge1\).
Average \(\#(\Sigma_j(\mathcal O_{F_j})\cap(\tau+[-R/2,R/2]^{2^j}))\) as \(\tau\) ranges over a fundamental domain for \(\Sigma_j(\mathcal O_{F_j})\).
By \eqref{eq:covolume-bound}, the average is
\[
\frac{\operatorname{vol}([-R/2,R/2]^{2^j})}{\operatorname{covol}(\Sigma_j(\mathcal O_{F_j}))}
=\frac{R^{2^j}}{|\Delta_{F_j}|^{1/2}}
\ge\left(\frac{R}{\sqrt D}\right)^{2^j}.
\]
Thus there exists \(\tau_{j,R}\in\mathbb R^{2^j}\) such that
\[
\#\bigl(\Sigma_j(\mathcal O_{F_j})\cap(\tau_{j,R}+[-R/2,R/2]^{2^j})\bigr)
\ge
\left(\frac{R}{\sqrt D}\right)^{2^j}.
\]
Define
\begin{equation}\label{eq:B-def}
B_{j,R}
:=
\{a\in\mathcal O_{F_j}:\Sigma_j(a)\in\tau_{j,R}+[-R/2,R/2]^{2^j}\}.
\end{equation}
Then
\begin{equation}\label{eq:B-size}
\#B_{j,R}\ge\left(\frac{R}{\sqrt D}\right)^{2^j}.
\end{equation}
For every \(a,c\in B_{j,R}\), the points \(\Sigma_j(a)\) and \(\Sigma_j(c)\) lie in the same translate of \([-R/2,R/2]^{2^j}\), and hence
\begin{equation}\label{eq:B-diff}
|\sigma_{j,k}(a-c)|\le R,
\qquad
k=1,\ldots,2^j.
\end{equation}

\subsection{Projection to the plane}
\label{subsec:projection-to-plane}

For \(j\) and \(R\), define
\begin{equation}\label{eq:P-def}
P_{j,R}
=
\{(\sigma_{j,1}(a),\sigma_{j,1}(b)):a,b\in B_{j,R}\}
\subset\mathbb R^2.
\end{equation}
Since \(\sigma_{j,1}\) is injective, \eqref{eq:B-size} gives
\begin{equation}\label{eq:P-size}
\#P_{j,R}=(\#B_{j,R})^2
\ge
\left(\frac{R}{\sqrt D}\right)^{2^{j+1}}.
\end{equation}
Note that the map \(B_{j,R}^2\to P_{j,R}\), \((a,b)\mapsto(\sigma_{j,1}(a),\sigma_{j,1}(b))\), is a bijection.
Thus a subset \(S\subset P_{j,R}\) corresponds to its inverse image
\begin{equation}\label{eq:S-lift}
\{(a,b)\in B_{j,R}^2:(\sigma_{j,1}(a),\sigma_{j,1}(b))\in S\}
\subset B_{j,R}^2.
\end{equation}
In the proof of Theorem~\ref{thm:main}, we choose \(j\) and \(R\) suitably and use \(P_{j,R}\) as the ambient planar set.

\begin{proposition}\label{prop:distance-sidon-lift}
Let \(S\subset P_{j,R}\), and suppose that all pairwise distances in \(S\) are distinct.
Let \(A\) be the inverse image of \(S\), as in \eqref{eq:S-lift}.
Then, for every nonzero \(\eta\in\mathcal O_{F_j}\),
\begin{equation}\label{eq:sidon-bound}
\#\{((a,b),(c,d))\in A\times A:
(a-c)^2+(b-d)^2=\eta\}
\le2.
\end{equation}
\end{proposition}

\begin{proof}
Let \((a,b)\) and \((c,d)\) be elements of \(A\).
The squared Euclidean distance between their projections is
\begin{equation}\label{eq:distance-lift}
(\sigma_{j,1}(a)-\sigma_{j,1}(c))^2
+(\sigma_{j,1}(b)-\sigma_{j,1}(d))^2
=
\sigma_{j,1}\bigl((a-c)^2+(b-d)^2\bigr).
\end{equation}
If \((a-c)^2+(b-d)^2=\eta\), then this squared distance is \(\sigma_{j,1}(\eta)\).
Since \(\eta\ne0\), the two projected endpoints are distinct.
Since all pairwise distances in \(S\) are distinct, there is at most one unordered pair at this distance, hence at most two ordered pairs.
\end{proof}

\section{The sieve estimate}
\label{sec:sieve-estimate}

We bound the size of a subset of \(P_{j,R}\) whose pairwise distances are all distinct.

\subsection{Counting elements in an ideal}
\label{subsec:counting-elements-in-an-ideal}

For a nonzero ideal \(\mathfrak a\subset\mathcal O_{F_j}\), write \(N\mathfrak a=\#(\mathcal O_{F_j}/\mathfrak a)\).\footnote{See, e.g., \cite[Section~4, ``Norms of ideals,'' before Proposition~4.2]{Mil20}.}
For \(\alpha\in F_j\), using the embeddings \(\sigma_{j,k}:F_j\hookrightarrow\mathbb R\) from \eqref{eq:real-embeddings}, write \(N_{F_j/\mathbb Q}(\alpha)=\prod_{k=1}^{2^j}\sigma_{j,k}(\alpha)\).\footnote{See, e.g., \cite[Section~2, ``Review of norms and traces,'' especially Corollary~2.20]{Mil20}.}
Recall from \eqref{eq:sigma-map} that \(\Sigma_j(\alpha)=(\sigma_{j,1}(\alpha),\ldots,\sigma_{j,2^j}(\alpha))\).

\begin{lemma}\label{lem:ideal-box}
Fix \(j\ge0\), a nonzero ideal \(\mathfrak a\subset\mathcal O_{F_j}\), and \(Y>0\).
Then
\[
\#\{\eta\in\mathfrak a:\Sigma_j(\eta)\in[-Y,Y]^{2^j}\}
\le
\left(1+\frac{2Y}{(N\mathfrak a)^{1/2^j}}\right)^{2^j}.
\]
\end{lemma}

\begin{proof}
If \(\eta_1\ne\eta_2\) are two elements of \(\mathfrak a\), then \(\eta_1-\eta_2\) is a nonzero element of \(\mathfrak a\).
Since \((\eta_1-\eta_2)\subset\mathfrak a\), we have
\[
|N_{F_j/\mathbb Q}(\eta_1-\eta_2)|
=
N((\eta_1-\eta_2))
\ge
N\mathfrak a.
\]
On the other hand,
\[
|N_{F_j/\mathbb Q}(\eta_1-\eta_2)|
=
\prod_{k=1}^{2^j} |\sigma_{j,k}(\eta_1-\eta_2)|.
\]
Since this product has \(2^j\) factors and is at least \(N\mathfrak a\), for some \(k\) we have
\[
|\sigma_{j,k}(\eta_1)-\sigma_{j,k}(\eta_2)|
=
|\sigma_{j,k}(\eta_1-\eta_2)|
\ge
(N\mathfrak a)^{1/2^j}.
\]
Therefore the points \(\Sigma_j(\eta)\), with \(\eta\in\mathfrak a\), are separated in the sup norm by \((N\mathfrak a)^{1/2^j}\).
This proves the lemma.
\end{proof}

\subsection{The sieve argument}
\label{subsec:sieve-argument}

Let \(t\ge1\) be an integer.
By Proposition~\ref{prop:tower}, we can choose distinct rational primes \(q_1,\ldots,q_t\) such that each \(q_i\equiv1\pmod4\) and each \(q_i\) splits completely in every \(F_j\).
Put
\[
Q=q_1\cdots q_t.
\]
Since \(q_i\) splits completely in \(F_j\),
\[
q_i\mathcal O_{F_j}
=
\mathfrak q_{i,j,1}\cdots\mathfrak q_{i,j,2^j},
\]
where \(\mathcal O_{F_j}/\mathfrak q_{i,j,k}\cong\mathbb F_{q_i}\) for \(k=1,\ldots,2^j\).
(See Section~\ref{subsec:basic-algebraic-number-theory}.)
Since \(q_i\equiv1\pmod4\), choose \(x_i\in\mathbb F_{q_i}\) with \(x_i^2=-1\).
For each \(k\), the inclusion \(\mathbb Z\subset\mathcal O_{F_j}\) induces an isomorphism \(\mathbb F_{q_i}=\mathbb Z/q_i\mathbb Z\to\mathcal O_{F_j}/\mathfrak q_{i,j,k}\), and we use \(x_i\) for the image of \(x_i\) under this isomorphism.

For each sign vector
\[
\theta=(\theta_{i,k})_{1\le i\le t,\ 1\le k\le 2^j}
\in\{\pm1\}^{t\,2^j},
\]
define the additive subgroup
\begin{equation}\label{eq:L-def}
L_{\theta,j}
=
\{(u,v)\in\mathcal O_{F_j}^2:
u\equiv\theta_{i,k}x_i v\pmod{\mathfrak q_{i,j,k}}
\text{ for every }i,k\}.
\end{equation}

\begin{lemma}\label{lem:isotropic-subgroups}
For every sign vector \(\theta\), we have \([\mathcal O_{F_j}^2:L_{\theta,j}]=Q^{2^j}\).
Moreover, if \((u,v)\in L_{\theta,j}\), then \(Q\mathcal O_{F_j}\mid u^2+v^2\).
\end{lemma}

\begin{proof}
For each prime ideal \(\mathfrak q_{i,j,k}\), the congruence
\[
u\equiv\theta_{i,k}x_i v\pmod{\mathfrak q_{i,j,k}}
\]
is the kernel of a surjective additive homomorphism
\[
(\mathcal O_{F_j}/\mathfrak q_{i,j,k})^2\to\mathcal O_{F_j}/\mathfrak q_{i,j,k},
\qquad
(u,v)\mapsto u-\theta_{i,k}x_i v.
\]
It therefore has index \(q_i\) modulo \(\mathfrak q_{i,j,k}\).
The Chinese remainder theorem over all pairs \((i,k)\) gives
\[
[\mathcal O_{F_j}^2:L_{\theta,j}]
=
\prod_{i=1}^t\prod_{k=1}^{2^j} q_i
=Q^{2^j}.
\]
If \((u,v)\in L_{\theta,j}\), then for every \(i,k\),
\[
u^2+v^2
\equiv
(\theta_{i,k}^2x_i^2+1)v^2
=0
\pmod{\mathfrak q_{i,j,k}}.
\]
Since \(Q=q_1\cdots q_t\), we have
\[
Q\mathcal O_{F_j}
=
\prod_{i=1}^t\prod_{k=1}^{2^j}\mathfrak q_{i,j,k}.
\]
Since the prime ideals \(\mathfrak q_{i,j,k}\) are distinct, their intersection is their product.
Hence \(Q\mathcal O_{F_j}\) divides \(u^2+v^2\).
\end{proof}

\begin{proposition}\label{prop:sieve}
For \(t\) sufficiently large, there exists a constant \(0<\rho<1\) such that the following holds.
If \(R\ge Q\) and \(S\subset P_{j,R}\) has all pairwise distances distinct, then
\begin{equation}\label{eq:sieve-conclusion}
\#S\le Q^{2^j}+\left(\frac{R}{\sqrt D}\right)^{2^j}\rho^{2^j}.
\end{equation}
\end{proposition}

\begin{proof}
Let \(A\subset B_{j,R}^2\) be the inverse image of \(S\) under the bijection \(B_{j,R}^2\to P_{j,R}\) from \eqref{eq:S-lift}.
Put \(\Lambda=\prod_{i=1}^t(1+1/q_i)\).
Choose \(t\) large enough so that
\begin{equation}\label{eq:rho-def}
\rho:=
\sqrt D\left(\frac{10\Lambda}{2^t}\right)^{1/2}
<1.
\end{equation}
This is possible because every \(q_i\ge5\), so \(\Lambda\le(6/5)^t\), and therefore \(10D\Lambda/2^t\le 10D(3/5)^t\to0\).
By \eqref{eq:rho-def},
\[
R^{2^j}\left(\frac{10\Lambda}{2^t}\right)^{2^j/2}
=
\left(\frac{R}{\sqrt D}\right)^{2^j}\rho^{2^j}.
\]
Thus, since \(\#S=\#A\), \eqref{eq:sieve-conclusion} follows from
\begin{equation}\label{eq:sieve-preliminary}
\#A\le Q^{2^j}+R^{2^j}\left(\frac{10\Lambda}{2^t}\right)^{2^j/2}.
\end{equation}
It remains to prove \eqref{eq:sieve-preliminary}.
If \(\#A<Q^{2^j}\), then \eqref{eq:sieve-preliminary} is immediate.
Thus assume from now on that \(\#A\ge Q^{2^j}\).

By Lemma~\ref{lem:isotropic-subgroups}, \([\mathcal O_{F_j}^2:L_{\theta,j}]=Q^{2^j}\).
If \(\#A\) is much larger than \(Q^{2^j}\), then there must be many \((p,p')\in A\times A\) with \(p\ne p'\) and \(p-p'\in L_{\theta,j}\).
Accordingly, we count triples
\begin{equation}\label{eq:T-def}
T=
\#\{(p,p',\theta)\in A\times A\times\{\pm1\}^{t2^j}:
p\ne p',\ p-p'\in L_{\theta,j}\}.
\end{equation}
For a fixed sign vector \(\theta\), let \(C_1,\ldots,C_{Q^{2^j}}\) be the \(L_{\theta,j}\)-cosets in \(\mathcal O_{F_j}^2\).
Then
\[
A=\bigsqcup_{s=1}^{Q^{2^j}} (A\cap C_s).
\]
Put \(n_s=\#(A\cap C_s)\).
Then \(\sum_{s=1}^{Q^{2^j}} n_s=\#A\), and the number of \((p,p')\in A\times A\) with \(p\ne p'\) and \(p-p'\in L_{\theta,j}\) is
\[
\sum_{s=1}^{Q^{2^j}} n_s(n_s-1)
=
\sum_{s=1}^{Q^{2^j}} n_s^2-\#A
\ge
\frac{(\#A)^2}{Q^{2^j}}-\#A
=
\frac{\#A(\#A-Q^{2^j})}{Q^{2^j}}.
\]
Summing over the \(2^{t2^j}\) sign vectors gives
\begin{equation}\label{eq:T-lower}
T\ge
2^{t2^j}\frac{\#A(\#A-Q^{2^j})}{Q^{2^j}}.
\end{equation}

We next prove an upper bound for \(T\).
For \((u,v)\in\mathcal O_{F_j}^2\), set
\begin{align*}
r(u,v)&=\#\{(p,p')\in A\times A:p-p'=(u,v)\},\\
w(u,v)&=\#\{\theta\in\{\pm1\}^{t2^j}:(u,v)\in L_{\theta,j}\}.
\end{align*}
By \eqref{eq:T-def},
\begin{equation}\label{eq:T-rw}
T=\sum_{(u,v)\ne(0,0)} r(u,v)w(u,v).
\end{equation}

By \eqref{eq:L-def},
\[
w(u,v)
=
\prod_{i=1}^t\prod_{k=1}^{2^j}w_{i,j,k}(u,v),
\]
where
\[
w_{i,j,k}(u,v)
=
\#\{\theta\in\{\pm1\}:u\equiv\theta x_i v\pmod{\mathfrak q_{i,j,k}}\}.
\]
Suppose \(w(u,v)>0\).
Let
\[
E=
\{(i,k):1\le i\le t,\ 1\le k\le 2^j,\ w_{i,j,k}(u,v)=2\}.
\]
Then \(w(u,v)=2^{\#E}\).
Since \(w(u,v)>0\), every \(\mathfrak q_{i,j,k}\) has \(w_{i,j,k}(u,v)\ge1\), and hence \(\mathfrak q_{i,j,k}\mid u^2+v^2\).
If \((i,k)\in E\), then
\[
u\equiv x_i v\pmod{\mathfrak q_{i,j,k}},
\qquad
u\equiv -x_i v\pmod{\mathfrak q_{i,j,k}}.
\]
Subtracting gives \(2x_i v\equiv0\pmod{\mathfrak q_{i,j,k}}\).
Since \(2x_i\ne0\) in \(\mathcal O_{F_j}/\mathfrak q_{i,j,k}\cong\mathbb F_{q_i}\), this gives \(v\equiv0\pmod{\mathfrak q_{i,j,k}}\), and then \(u\equiv0\pmod{\mathfrak q_{i,j,k}}\).
Thus \(\mathfrak q_{i,j,k}^2\mid u^2+v^2\) for every \((i,k)\in E\).
For every subset \(E'\subset E\), put
\[
\mathfrak b=\prod_{(i,k)\in E'}\mathfrak q_{i,j,k}.
\]
Then \(\mathfrak b\mid Q\mathcal O_{F_j}\) and \(Q\mathfrak b\mid u^2+v^2\).
Therefore we obtain
\begin{equation}\label{eq:w-bound}
w(u,v)
\le
2^{\#E}
\le
\#\{\mathfrak b\mid Q\mathcal O_{F_j} : Q\mathfrak b\mid u^2+v^2\}.
\end{equation}
The bound \eqref{eq:w-bound} is immediate if \(w(u,v)=0\).

By Proposition~\ref{prop:distance-sidon-lift}, for every nonzero \(\eta\in\mathcal O_{F_j}\),
\begin{equation}\label{eq:sidon-bound-difference}
\sum_{u^2+v^2=\eta} r(u,v)\le2.
\end{equation}
Also, since \(F_j\) is totally real, \(u^2+v^2=0\) implies \(u=v=0\); hence \(\eta=0\) contributes nothing to \(T\), where \((u,v)\ne(0,0)\).

For \((u,v)\in A-A\), \eqref{eq:B-diff} gives
\[
|\sigma_{j,k}(u)|\le R,
\qquad
|\sigma_{j,k}(v)|\le R,
\qquad
k=1,\ldots,2^j.
\]
Therefore
\begin{equation}\label{eq:eta-box}
\Sigma_j(u^2+v^2)\in[-2R^2,2R^2]^{2^j}.
\end{equation}
Combining \eqref{eq:T-rw}, \eqref{eq:w-bound}, \eqref{eq:sidon-bound-difference}, and \eqref{eq:eta-box} gives
\begin{align}
T
&=
\sum_{(u,v)\ne(0,0)} r(u,v)w(u,v)
\notag\\
&\le
\sum_{(u,v)\ne(0,0)} r(u,v)
\sum_{\mathfrak b\mid Q\mathcal O_{F_j},\ Q\mathfrak b\mid u^2+v^2}1
\notag\\
&\le
\sum_{\mathfrak b\mid Q\mathcal O_{F_j}}
\sum_{\substack{0\ne\eta\in Q\mathfrak b\\
\Sigma_j(\eta)\in[-2R^2,2R^2]^{2^j}}}
\sum_{u^2+v^2=\eta}r(u,v)
\notag\\
&\le
2\sum_{\mathfrak b\mid Q\mathcal O_{F_j}}
\#\{\eta\in Q\mathfrak b:\Sigma_j(\eta)\in[-2R^2,2R^2]^{2^j}\}.
\label{eq:T-upper-prebox}
\end{align}

By multiplicativity of the ideal norm, the factorization
\[
Q\mathcal O_{F_j}
=
\prod_{i=1}^t\prod_{k=1}^{2^j}\mathfrak q_{i,j,k}
\]
and the equality \(N\mathfrak q_{i,j,k}=\#(\mathcal O_{F_j}/\mathfrak q_{i,j,k})=\#\mathbb F_{q_i}=q_i\) give
\[
N(Q\mathcal O_{F_j})
=
\prod_{i=1}^t\prod_{k=1}^{2^j}q_i
=
Q^{2^j}.
\]
Thus
\begin{equation}\label{eq:N-Qb}
(N(Q\mathfrak b))^{1/2^j}
=
(N(Q\mathcal O_{F_j})N\mathfrak b)^{1/2^j}
=
Q(N\mathfrak b)^{1/2^j}.
\end{equation}
Also, since \(\mathfrak b\mid Q\mathcal O_{F_j}\), we have \((N\mathfrak b)^{1/2^j}\le Q\).
Because \(R\ge Q\),
\begin{equation}\label{eq:R-b-lower}
\frac{R^2}{Q(N\mathfrak b)^{1/2^j}}\ge1.
\end{equation}
Applying Lemma~\ref{lem:ideal-box} to \(\mathfrak a=Q\mathfrak b\) and \(Y=2R^2\), and using \eqref{eq:N-Qb} and \eqref{eq:R-b-lower}, we get
\[
\#\{\eta\in Q\mathfrak b:\Sigma_j(\eta)\in[-2R^2,2R^2]^{2^j}\}
\le
\left(1+\frac{4R^2}{(N(Q\mathfrak b))^{1/2^j}}\right)^{2^j}
\le
\left(\frac{5R^2}{Q}\right)^{2^j}\frac1{N\mathfrak b}.
\]
Substituting this into \eqref{eq:T-upper-prebox} gives
\[
T
\le
\left(\frac{10R^2}{Q}\right)^{2^j}
\sum_{\mathfrak b\mid Q\mathcal O_{F_j}}\frac1{N\mathfrak b}.
\]
Because \(N\mathfrak q_{i,j,k}=q_i\) for every \(k\),
\[
\sum_{\mathfrak b\mid Q\mathcal O_{F_j}}\frac1{N\mathfrak b}
=
\prod_{i=1}^t\prod_{k=1}^{2^j}
\left(1+\frac1{N\mathfrak q_{i,j,k}}\right)
=
\prod_{i=1}^t\prod_{k=1}^{2^j}
\left(1+\frac1{q_i}\right)
=
\Lambda^{2^j}.
\]
Hence
\begin{equation}\label{eq:T-upper}
T\le
\left(\frac{10R^2\Lambda}{Q}\right)^{2^j}.
\end{equation}

Comparing \eqref{eq:T-lower} and \eqref{eq:T-upper}, we obtain
\[
2^{t2^j}\frac{\#A(\#A-Q^{2^j})}{Q^{2^j}}
\le
\left(\frac{10R^2\Lambda}{Q}\right)^{2^j}.
\]
Therefore
\[
\#A(\#A-Q^{2^j})
\le
\left(\frac{10R^2\Lambda}{2^t}\right)^{2^j}.
\]
Since \(\#A\ge Q^{2^j}\), we have \((\#A-Q^{2^j})^2\le \#A(\#A-Q^{2^j})\), and hence
\[
\#A-Q^{2^j}
\le
R^{2^j}\left(\frac{10\Lambda}{2^t}\right)^{2^j/2}.
\]
This proves \eqref{eq:sieve-preliminary}.
\end{proof}

\section{Proof of Theorem~\ref{thm:main}}
\label{sec:proof-main-theorem}

Choose \(t\) large enough so that Proposition~\ref{prop:sieve} applies.
We use the corresponding primes \(q_1,\ldots,q_t\), the product \(Q=q_1\cdots q_t\), and the constant \(0<\rho<1\) from Section~\ref{subsec:sieve-argument}.
Choose a real number \(a>\log Q\).

Let \(n\) be sufficiently large.
Let \(j\ge0\) be the integer satisfying
\begin{equation}\label{eq:j-choice}
\frac{\log n}{4a}<2^j\le\frac{\log n}{2a}.
\end{equation}
Put \(R=\sqrt D\,n^{1/2^{j+1}}\).
The upper bound in \eqref{eq:j-choice} gives
\[
\log\left(\frac{R}{\sqrt D}\right)
=
\frac{\log n}{2^{j+1}}
\ge a,
\]
so \(R\ge\sqrt D\,e^a\ge Q\).

Let \(B_{j,R}\) and \(P_{j,R}\) be as in \eqref{eq:B-def} and \eqref{eq:P-def}.
By \eqref{eq:P-size},
\[
\#P_{j,R}
\ge
\left(\frac{R}{\sqrt D}\right)^{2^{j+1}}
=n.
\]
Let \(P_n\subset P_{j,R}\) be any \(n\)-point subset.

Let \(S\subset P_n\) have all pairwise distances distinct.
Since \(S\subset P_{j,R}\), Proposition~\ref{prop:sieve} gives
\[
\#S
\le
Q^{2^j}+
\left(\frac{R}{\sqrt D}\right)^{2^j}\rho^{2^j}
=
Q^{2^j}+n^{1/2}\rho^{2^j}.
\]
By \eqref{eq:j-choice},
\[
Q^{2^j}
=
\exp(2^j\log Q)
\le
\exp\left(\frac{\log Q}{2a}\log n\right)
=
n^{\log Q/(2a)}
\]
and
\[
n^{1/2}\rho^{2^j}
=
n^{1/2}\exp(2^j\log\rho)
\le
n^{1/2}\exp\left(\frac{\log\rho}{4a}\log n\right)
=
n^{1/2+\log\rho/(4a)},
\]
we may choose \(\varepsilon>0\) such that
\[
\frac{\log Q}{2a}\le\frac12-\varepsilon,
\qquad
\frac12+\frac{\log\rho}{4a}\le\frac12-\varepsilon.
\]
Then
\[
\#S\le2n^{1/2-\varepsilon}.
\]
Thus every subset of \(P_n\) with all pairwise distances distinct has size at most \(2n^{1/2-\varepsilon}\).
Consequently
\[
F_2(n)
\le
2n^{1/2-\varepsilon}
\]
for all sufficiently large \(n\).

\section*{Use of AI}

During the preparation of this work, the author used OpenAI's GPT-5.5 Pro to explore possible proof strategies and to obtain methodological ideas behind the argument.
After using this tool, the author independently verified the mathematical details, refined the exposition, and takes full responsibility for the final manuscript.


\bibliographystyle{amsplain}
\bibliography{references}

\end{document}
